% ============================================================================
% Fichier  : partie2/P02_C04_rgpd.tex
% Titre    : Le RGPD et son Effet hors d'Europe
% Auteur   : MINKA MI NGUIDJOÏ Thierry Emmanuel
% Version  : 3.0 -- novembre 2026
% Corrigés : corriges/Corr_C04.tex
% ============================================================================
% Sources : règlement (UE) 2016/679 (Journal officiel de l'UE), paraphrasé ;
% lignes directrices 3/2018 du Comité européen de la protection des données
% sur le champ territorial ; décision d'exécution (UE) 2021/914 (clauses
% contractuelles types). Loi n° 2024/017, texte officiel. Le service NOUN
% DIASPORA et le contrat bruxellois de MEDISOFT sont des ajouts fictifs de
% la version 3.
% ============================================================================

\chapter{Le RGPD et son Effet hors d'Europe}
\label{chap:rgpd}
\chaptermark{Le RGPD hors d'Europe}

\epigraph{%
    « Le règlement européen ne s'arrête pas aux frontières de l'Europe.\\
    Il ne remplace pas pour autant la loi du pays où l'on est. »%
}{Principe du cours}

\niveauChapitre{Fondamental}{%
    Chapitres~\ref{chap:principes} et~\ref{chap:international}
}

\begin{boxobjectifs}
\begin{itemize}
    \item Présenter l'économie du RGPD en une quinzaine d'articles utiles à
          l'auditeur.
    \item Dire quand le RGPD oblige une organisation établie au Cameroun, en
          appliquant son article~3.
    \item Distinguer les trois voies d'exposition~: directe, contractuelle,
          en cascade.
    \item Comparer la loi n°~2024/017 et le RGPD point par point, sans
          prêter à l'une le contenu de l'autre.
    \item Coter une même exigence contre les deux textes quand les deux
          s'appliquent.
\end{itemize}
\end{boxobjectifs}

\bigskip

% ============================================================================
\section*{Avant de commencer : « c'est le RGPD camerounais »}
% ============================================================================

Bafoussam, lundi 3~août~2026. Le directeur général de NYANGA CRÉDIT a reçu
d'un cabinet de conseil une offre de « mise en conformité RGPD camerounais »
en douze semaines. Le document promet la désignation d'un délégué à la
protection des données « obligatoire », une procédure de notification des
violations « sous 72~heures » et la constitution d'un dossier d'« intérêt
légitime » pour la vidéosurveillance.

M.~NGANDO pose la proposition à côté du texte de la loi. Aucune de ces trois
affirmations ne s'y trouve. La loi ne connaît pas de délégué obligatoire,
elle dit « sans délai » et non 72~heures et elle ignore l'intérêt légitime.
L'offre ne dit pas un mot, en revanche, de l'autorisation préalable de
l'article~19.

Le même jour, à quelques rues de là, M.~NJOYA pose à Mme~BELLO la question
inverse. NOUN PAY\footnote{NOUN PAY, son service NOUN DIASPORA, MEDISOFT
SARL et la clinique bruxelloise citée plus bas sont fictifs.} a ouvert en
mars un service de transfert d'argent pour les Camerounais de France et de
Belgique. « Eux sont en Europe. Est-ce que leur règlement nous concerne~? »
La réponse est oui.

\begin{boxalert}
\textbf{Violation typique.} Deux erreurs symétriques. Lire la loi
camerounaise à travers le RGPD et lui attribuer des obligations européennes
qu'elle ne contient pas. Ou croire qu'une entreprise camerounaise échappe au
RGPD parce qu'elle n'a aucun bureau en Europe.

\textbf{Règle~:} l'auditeur détermine d'abord, traitement par traitement,
quel texte s'applique. La loi n°~2024/017 s'applique à tout traitement visé
par son article~2. Le RGPD s'y ajoute quand son article~3 est rempli ou
qu'un contrat l'impose. Aucun des deux ne dispense de l'autre.
\end{boxalert}

% ============================================================================
\section{Le RGPD en Quinze Articles}
\label{sec:c04-economie}
% ============================================================================

Le règlement (UE) 2016/679~\cite{ue2016rgpd} est applicable depuis le
25~mai~2018 dans tous les États de l'Union européenne. Il compte
quatre-vingt-dix-neuf articles. L'auditeur en retient une quinzaine.

\begin{tableaucours}{cL{4.0cm}L{6.6cm}}{%
    \tableauHeader{Art.} &
    \tableauHeader{Objet} &
    \tableauHeader{Ce qu'il faut en retenir (paraphrasé)}
}
3 & Champ territorial & Établissement dans l'Union, ou ciblage de personnes
qui s'y trouvent \\
5 & Principes & Sept principes, dont la minimisation et la responsabilité \\
6 & Bases légales & Six bases, dont le contrat et l'intérêt légitime \\
9 & Catégories particulières & Liste fermée de données sensibles \\
12 à 22 & Droits & Information, accès, rectification, effacement,
limitation, portabilité, opposition~; réponse en un mois \\
24, 25 & Responsabilité & Démontrer la conformité~; protection dès la
conception \\
27 & Représentant & Obligatoire pour l'acteur hors Union visé par
l'article~3 \\
28 & Sous-traitant & Contrat écrit au contenu imposé \\
30 & Registre & Tenu par le responsable et par le sous-traitant \\
32 & Sécurité & Mesures adaptées au risque \\
33, 34 & Violations & Autorité en 72~heures~; personnes si le risque est
élevé \\
35, 36 & Analyse d'impact & Obligatoire si le risque est élevé~;
consultation de l'autorité \\
37 à 39 & Délégué & Obligatoire dans trois cas \\
44 à 49 & Transferts & Adéquation, garanties appropriées, dérogations \\
83 & Amendes & Jusqu'à 20~millions d'euros ou 4~\% du chiffre d'affaires
mondial \\
\end{tableaucours}
\vspace{-0.8em}
\begin{center}{\small\itshape Tableau~4.1. Quinze articles du RGPD pour
l'auditeur}\end{center}

Le règlement repose sur une idée que la loi camerounaise n'a pas retenue~:
la \termecle{responsabilité}. Il a supprimé les formalités préalables auprès
de l'autorité. En échange, le responsable doit pouvoir démontrer à tout
moment qu'il respecte le texte. La loi n°~2024/017 fait le choix opposé,
celui du contrôle préalable par l'autorisation (Chapitre~\ref{chap:loi}).

% ============================================================================
\section{Quand le RGPD Oblige une Organisation Camerounaise}
\label{sec:c04-article3}
% ============================================================================

L'article~3 du règlement fixe son champ territorial. Il retient deux
critères, dont le second vise précisément les acteurs établis hors de
l'Union.

\begin{enumerate}
    \item \textbf{L'établissement.} Le règlement s'applique aux traitements
          faits dans le cadre des activités d'un établissement situé dans
          l'Union, que le traitement ait lieu ou non dans l'Union. Une
          filiale, une succursale ou un simple bureau stable suffisent.
    \item \textbf{Le ciblage.} Il s'applique aussi à un responsable ou à un
          sous-traitant non établi dans l'Union, lorsque le traitement est
          lié à l'offre de biens ou de services à des personnes qui se
          trouvent dans l'Union, ou au suivi de leur comportement dans
          l'Union.
\end{enumerate}

Le critère du ciblage tient à la localisation de la personne au moment du
traitement, pas à sa nationalité. Un Camerounais qui vit à Lyon est protégé
par le RGPD. Un Français en visite à Kribi ne l'est pas pour les traitements
d'une entreprise camerounaise qui ne vise pas l'Europe~; il l'est en
revanche par la loi camerounaise, dont l'article~2 protège les personnes
« en transit ».

Le Comité européen de la protection des données a précisé ce qu'est une
offre~\cite{edpb2018territorial}. Un site simplement accessible depuis
l'Europe ne suffit pas. Il faut des indices d'une intention de s'adresser à
des personnes dans l'Union.

\begin{tableaucours}{L{6.0cm}L{6.0cm}}{%
    \tableauHeader{Indice de ciblage} &
    \tableauHeader{À NOUN DIASPORA}
}
Prix affichés dans une monnaie de l'Union & Frais affichés en euros \\
Langue ou pays cités dans l'offre & « Envoyez depuis la France et la
Belgique » \\
Publicité achetée pour un public de l'Union & Campagnes sur les réseaux
sociaux, zone France et Belgique \\
Moyens de paiement de l'Union acceptés & Cartes bancaires européennes \\
Mention de clients de l'Union & Témoignages de clients de Lyon et de Liège \\
\end{tableaucours}
\vspace{-0.8em}
\begin{center}{\small\itshape Tableau~4.2. Les indices de ciblage, appliqués
au service de NOUN PAY}\end{center}

Les cinq indices sont réunis. NOUN PAY est un responsable de traitement non
établi dans l'Union qui offre un service à des personnes qui s'y trouvent.
Le règlement s'applique aux traitements de NOUN DIASPORA et à eux seuls~:
le microcrédit proposé aux clients de Bafoussam n'en relève pas.

L'article~27 en tire une première conséquence. L'acteur ainsi visé désigne
par écrit un représentant établi dans l'un des États où se trouvent les
personnes. Le représentant reçoit les demandes des personnes et de
l'autorité.

% ============================================================================
\section{Trois Voies d'Exposition}
\label{sec:c04-voies}
% ============================================================================

\begin{tableaucours}{L{2.4cm}L{4.6cm}L{5.0cm}}{%
    \tableauHeader{Voie} &
    \tableauHeader{Situation} &
    \tableauHeader{Conséquence}
}
Directe & L'organisation camerounaise cible des personnes dans l'Union
(NOUN DIASPORA) & Le RGPD lui est opposable pour ces traitements~;
représentant à désigner \\
Contractuelle & Elle est sous-traitante d'un responsable européen & Contrat
de l'article~28 et clauses contractuelles types pour le transfert vers le
Cameroun \\
En cascade & Elle confie ses données à un prestataire établi dans l'Union &
Le prestataire applique ses obligations de sous-traitant et les répercute
dans son contrat \\
\end{tableaucours}
\vspace{-0.8em}
\begin{center}{\small\itshape Tableau~4.3. Trois manières de rencontrer le
RGPD depuis le Cameroun}\end{center}

\subsection{La voie contractuelle}

MEDISOFT SARL, l'éditeur de Douala qui héberge le dossier patient de MBOLO,
a signé en 2026 un contrat de télémaintenance avec une clinique de
Bruxelles. Ses techniciens consultent depuis Douala des données de patients
belges. La clinique est responsable de traitement, MEDISOFT sous-traitant.

La clinique ne peut pas transférer ces données au Cameroun sans garantie~:
aucune décision d'adéquation ne couvre le pays
(Chapitre~\ref{chap:international}). Elle recourt donc aux clauses
contractuelles types adoptées par la Commission européenne en
2021~\cite{ue2021clauses}. MEDISOFT les signe telles quelles, car elles ne
se négocient pas. Elles lui imposent de n'agir que sur instruction,
d'assurer la sécurité, de signaler toute violation, d'accepter des audits et
de dire si le droit camerounais permet à une autorité d'accéder aux données.

Le RGPD est ici \textbf{contractuel} au sens du
Chapitre~\ref{chap:international}. L'auditeur de MEDISOFT cote ces
engagements si le contrat figure dans ses critères.

\subsection{La voie en cascade}

L'hébergeur de NOUN PAY est établi dans l'Union. Il est soumis au RGPD pour
ses propres obligations de sous-traitant, quel que soit le pays de son
client. Son contrat type reprend donc les clauses de l'article~28 et NOUN
PAY les a acceptées en ouvrant son compte. L'effet est favorable~: NOUN PAY
dispose d'un droit d'audit, d'un engagement de notification et d'une liste
des sous-traitants ultérieurs qu'elle n'aurait peut-être pas obtenus seule.

Un point reste à surveiller. Quand l'hébergeur renvoie les données vers le
Cameroun, par exemple pour qu'un agent de Bafoussam les consulte, il fait
lui-même un transfert hors de l'Union. Il peut exiger pour cela des
garanties de son client.

% ============================================================================
\section{La Loi et le RGPD, Point par Point}
\label{sec:c04-comparaison}
% ============================================================================

\begin{tableaucours}{L{2.8cm}L{4.6cm}L{4.6cm}}{%
    \tableauHeader{Point} &
    \tableauHeader{Loi n°~2024/017} &
    \tableauHeader{RGPD}
}
Champ & Personnes établies, résidant ou en transit au Cameroun~; acteurs
établis au Cameroun (art.~2) & Établissement dans l'Union~; ciblage de
personnes dans l'Union (art.~3) \\
Formalité préalable & Autorisation pour tout traitement (art.~19) & Aucune,
sauf consultation préalable (art.~36) \\
Bases du traitement & Consentement et trois dérogations (art.~9) & Six
bases, dont contrat et intérêt légitime (art.~6) \\
Minimisation & Pas de principe général & Principe (art.~5) \\
Délégué & Non prévu & Obligatoire dans trois cas (art.~37) \\
Registre & Responsable et sous-traitant (art.~29) & Les deux, avec une
exemption limitée pour les petites structures (art.~30) \\
Évaluation d'impact & Si risque élevé~; modalités réglementaires (art.~33) &
Si risque élevé~; critères publiés (art.~35) \\
Violation & « Sans délai », à l'Autorité et à la personne, pour toute
violation (art.~22) & 72~heures à l'autorité~; personnes si risque élevé
(art.~33, 34) \\
Sous-traitant & Mêmes obligations que le responsable (art.~20) & Obligations
propres, plus limitées (art.~28) \\
Responsables conjoints & Convention rendue publique (art.~31) & Grandes
lignes mises à disposition (art.~26) \\
Transferts & Autorisation, quatre vérifications (art.~32) & Adéquation,
garanties, dérogations (art.~44 à~49) \\
Décision automatisée & Droit d'opposition~; profilage puni (art.~44, 65) &
Encadrée, non interdite (art.~22) \\
Données sensibles & Liste ouverte, dont transactions bancaires et origine
régionale (art.~5) & Liste fermée (art.~9) \\
Délai de réponse aux demandes & Renvoyé à un texte réglementaire (art.~47) &
Un mois, prolongeable (art.~12) \\
Sanction maximale & Un milliard de francs CFA, personne morale (art.~71) &
4~\% du chiffre d'affaires mondial (art.~83) \\
\end{tableaucours}
\vspace{-0.8em}
\begin{center}{\small\itshape Tableau~4.4. Quinze points de comparaison}\end{center}

Trois lectures de ce tableau sont utiles.

La loi camerounaise est \textbf{plus stricte} sur six points~: l'autorisation
préalable, le petit nombre de fondements, la notification de toute
violation, les obligations du sous-traitant, la publicité de la convention
entre responsables conjoints et la liste des données sensibles. Une
organisation conforme au RGPD n'est donc pas conforme à la loi.

Elle est \textbf{moins précise} sur quatre~: la minimisation, les délais,
les critères du risque élevé et le contenu de la notification. Là, le RGPD
fournit des repères que l'auditeur peut emprunter.

Elle est \textbf{simplement différente} sur le reste. Le délégué n'existe
pas dans la loi et sa désignation ne peut être qu'une bonne pratique.

% ============================================================================
\section{Emprunter sans Confondre}
\label{sec:c04-emprunter}
% ============================================================================

Le Chapitre~\ref{chap:audit} montrera qu'un critère renvoyant à un texte non
publié se remplace par un \termecle{critère de substitution} déclaré. Le
droit européen est la première source de ces critères, parce qu'il est
précis, publié et commenté. Trois emprunts reviennent dans ce cours.

\begin{tableaucours}{L{3.6cm}L{4.2cm}L{4.2cm}}{%
    \tableauHeader{Silence de la loi} &
    \tableauHeader{Repère européen} &
    \tableauHeader{Formulation dans le rapport}
}
Délai de réponse aux demandes (art.~47) & Un mois (RGPD, art.~12) & « 30
jours (critère de l'audit) » \\
Critères du risque élevé (art.~33) & Liste du Comité européen, méthode de la
CNIL~\cite{cnil2018pia} & « Critères déclarés au plan d'audit » \\
Contenu de la notification (art.~22) & Rubriques des articles~33 et~34 &
« Contenu minimal retenu par l'audit » \\
\end{tableaucours}
\vspace{-0.8em}
\begin{center}{\small\itshape Tableau~4.5. Trois emprunts et leur
formulation}\end{center}

La règle tient en une phrase. Un emprunt se déclare dans le plan d'audit, se
nomme dans le rapport et ne se présente jamais comme une obligation légale
camerounaise. L'emprunt ne joue que dans un sens~: il précise un silence, il
n'allège jamais une exigence de la loi. On n'importe pas l'intérêt légitime.

% ============================================================================
\section{La Double Conformité}
\label{sec:c04-double}
% ============================================================================

Pour NOUN DIASPORA, les deux textes s'appliquent au même traitement. Le
client est à Lyon, donc protégé par le RGPD. NOUN PAY est établie au
Cameroun, donc soumise à la loi par son article~2. L'auditeur cote chaque
exigence deux fois.

\begin{tableaucours}{L{3.4cm}L{3.4cm}L{3.4cm}L{1.6cm}}{%
    \tableauHeader{Exigence} &
    \tableauHeader{Critère camerounais} &
    \tableauHeader{Critère européen} &
    \tableauHeader{Règle}
}
Fondement du traitement & Consentement ou dérogation (art.~9) & L'une des
six bases (art.~6) & Le plus strict \\
Notification d'une violation & Sans délai, Autorité et personne (art.~22) &
72~heures, autorité~; personne si risque élevé & Les deux \\
Formalité préalable & Autorisation (art.~19) & Aucune & Camerounais \\
Représentant & Aucun & Désignation écrite (art.~27) & Européen \\
\end{tableaucours}
\vspace{-0.8em}
\begin{center}{\small\itshape Tableau~4.6. Une exigence, deux critères}\end{center}

Quand les deux textes demandent la même chose à des degrés différents,
l'organisation satisfait le plus strict et satisfait ainsi les deux. Quand
ils demandent des choses différentes, elle fait les deux. Les deux textes ne
se contredisent sur aucun de ces points~: aucun n'interdit ce que l'autre
impose.

\begin{boxscenario}{Mission NOUN PAY. NOUN DIASPORA face au RGPD, septembre 2026}
\textbf{Constatations.}
\begin{itemize}
    \item 6\,840 clients inscrits depuis la France et la Belgique entre mars
          et août~2026~; frais en euros, publicité ciblée sur ces deux pays.
    \item Aucun représentant désigné dans l'Union. Aucune mention du RGPD
          dans la notice de l'application.
    \item La vérification d'identité de ces clients passe par le même
          prestataire de selfies, dont le pays reste inconnu.
    \item Le contrat de l'hébergeur contient les clauses de l'article~28.
    \item Aucune analyse d'impact~; aucune procédure de notification.
\end{itemize}

\textbf{Analyse} (confiance élevée). Le service remplit le critère du
ciblage. NOUN PAY l'a lancé sans se demander quel droit protégeait ses
clients. Le seul point conforme lui vient de son hébergeur.
\end{boxscenario}

\begin{boxlecture}{NOUN DIASPORA, exigences du RGPD}
\begin{tableaucours}{L{7.6cm}cL{3.4cm}}{%
    \tableauHeader{Exigence} &
    \tableauHeader{Cote} &
    \tableauHeader{Preuve}
}
Un représentant est désigné dans l'Union (art.~27) & \NON{} & Aucune
désignation \\
Chaque traitement a une base légale identifiée (art.~6) & \NON{} & Case
unique \\
L'information de l'article~13 est fournie & \NON{} & Notice de
l'application \\
Un registre est tenu (art.~30) & \NON{} & Aucun registre \\
Une analyse d'impact a été conduite (art.~35) & \NON{} & Aucune \\
Une procédure permet de notifier en 72~heures (art.~33) & \NON{} & Aucune \\
Le contrat de l'hébergeur contient les clauses de l'article~28 & \OUI{} &
Contrat, annexe données \\
Un délégué doit être désigné (art.~37) & \NEK{} & Suivi « à grande échelle »
à apprécier \\
\end{tableaucours}

\medskip
\textbf{Décompte.} Un \OUI{}, six \NON{}, un \NEK{} sur huit.
\hfill\lecture{0{,}125}{0{,}125}{1{,}06}

\medskip
\textbf{Lecture.} Le domaine est clair et non conforme. Ces huit cotes
s'ajoutent à celles de la loi camerounaise, elles ne les remplacent pas~: le
même service reçoit par ailleurs un \NON{} au titre de l'article~19.
\end{boxlecture}

% ============================================================================
\section{Ce que le RGPD demande à chaque Métier}
\label{sec:c04-metiers}
% ============================================================================

\begin{boxmetier}{Auditeur de la protection des données}
Poser la question de l'article~3 pour chaque traitement et consigner la
réponse au plan d'audit. Tenir deux colonnes de cotes quand les deux textes
s'appliquent.
\end{boxmetier}

\begin{boxmetier}{Juriste conformité}
Lire avant signature les clauses qu'un client ou un prestataire européen
impose. Elles créent des obligations d'audit, de notification et de
coopération que l'organisation devra tenir.
\end{boxmetier}

\begin{boxmetier}{Responsable de traitement}
Décider en connaissance de cause d'ouvrir un service à des personnes en
Europe. Le marché de la diaspora emporte le RGPD avec lui.
\end{boxmetier}

\begin{boxresume}
\begin{itemize}
    \item Le RGPD repose sur la responsabilité du responsable, la loi
          n°~2024/017 sur l'autorisation préalable.
    \item Son article~3 l'applique à une organisation camerounaise qui cible
          des personnes se trouvant dans l'Union~; l'article~27 lui impose
          alors un représentant.
    \item Trois voies d'exposition~: directe, contractuelle, en cascade.
    \item La loi camerounaise est plus stricte que le RGPD sur six points et
          moins précise sur quatre~; être conforme à l'un ne dit rien de
          l'autre.
    \item Un repère européen s'emprunte comme critère de substitution
          déclaré, pour préciser un silence, jamais pour alléger la loi.
\end{itemize}
\end{boxresume}

\begin{boxinfo}
\textbf{Points de contrôle.} Vous avez compris ce chapitre si vous savez~:
\begin{itemize}
    \item appliquer les indices de ciblage à un service en ligne~;
    \item citer trois points où la loi camerounaise est plus stricte que le
          RGPD~;
    \item formuler dans un rapport un critère emprunté au droit européen.
\end{itemize}
\end{boxinfo}

\begin{boxexo}[title={\ding{45}\quad Exercice~4.1. Le RGPD s'applique-t-il \niveaubase}]
Dites si le RGPD s'applique et par quelle voie~: (a) NYANGA CRÉDIT gère les
comptes de ses clients de l'Ouest~; (b) NOUN PAY propose NOUN DIASPORA~;
(c) MEDISOFT assure la télémaintenance d'une clinique de Bruxelles~;
(d) un touriste allemand est soigné à MBOLO~; (e) NOUN PAY héberge ses
données à Francfort.
\end{boxexo}

\begin{boxexo}[title={\ding{45}\quad Exercice~4.2. L'offre du cabinet \niveaumoyen}]
Reprenez l'offre de « mise en conformité RGPD camerounais » reçue par
NYANGA. Corrigez ses trois affirmations et ajoutez les trois obligations
camerounaises qu'elle omet, articles à l'appui.
\end{boxexo}

\begin{boxexo}[title={\ding{45}\quad Exercice~4.3. Deux colonnes \niveaumoyen}]
NOUN PAY découvre le 2~octobre qu'un fichier de 1\,200 clients de NOUN
DIASPORA a été envoyé par erreur à un partenaire. Dressez le tableau des
obligations de notification au titre de l'article~22 de la loi et des
articles~33 et~34 du RGPD~: destinataires, délais, contenu. Que fait NOUN
PAY pour satisfaire les deux textes~?
\end{boxexo}

\begin{boxexo}[title={\ding{45}\quad Exercice~4.4. Les clauses de MEDISOFT \niveauexpert}]
Les clauses signées par MEDISOFT l'obligent à dire si le droit camerounais
permet à une autorité publique d'accéder aux données des patients belges.
\begin{enumerate}[(a)]
    \item Quels textes camerounais MEDISOFT doit-il examiner pour répondre~?
          Appuyez-vous sur le Chapitre~\ref{chap:loi2010}.
    \item Le transfert de la Belgique vers le Cameroun est-il aussi un
          traitement soumis à l'article~19 de la loi n°~2024/017~?
    \item Comment l'auditeur de MEDISOFT cote-t-il la proposition « MEDISOFT
          respecte ses engagements envers la clinique »~?
\end{enumerate}
\end{boxexo}

\bigskip

\begin{boxinfo}
\textbf{Pour aller plus loin.}
\begin{itemize}
    \item Le RGPD~\cite{ue2016rgpd}, articles~3, 27, 28 et~44 à~49.
    \item Les lignes directrices 3/2018 du Comité européen sur le champ
          territorial~\cite{edpb2018territorial}.
    \item Les recommandations 01/2020 sur les transferts~\cite{edpb2020transferts}.
\end{itemize}
\end{boxinfo}

\begin{boxinfo}
\textbf{Transition.} Le RGPD est une loi. D'autres référentiels n'obligent
personne et organisent pourtant le travail de milliers d'entreprises~: le
cadre du NIST et les normes ISO. C'est l'objet du
Chapitre~\ref{chap:normes}.
\end{boxinfo}
